\documentclass[12pt, a4paper]{article}

\usepackage{geometry}
\geometry{a4paper, top=2.5cm, bottom=2.5cm, left=2.5cm, right=3cm}
\usepackage{graphicx}
\usepackage{booktabs}
\usepackage{longtable}
\usepackage{array}
\usepackage{multirow}
\usepackage{xcolor}
\usepackage{hyperref}
\usepackage{fancyhdr}
\usepackage{enumitem}
\usepackage{fontspec}
\usepackage{xepersian}
\settextfont[Scale=1.1]{Vazirmatn}
\setlatintextfont{Times New Roman}

\definecolor{primaryorange}{RGB}{234, 88, 12}
\definecolor{lightgray}{RGB}{245, 245, 245}
\definecolor{darkgray}{RGB}{80, 80, 80}
\definecolor{successgreen}{RGB}{22, 163, 74}
\definecolor{infoblue}{RGB}{37, 99, 235}

\hypersetup{
    colorlinks=true,
    linkcolor=primaryorange,
    urlcolor=primaryorange
}

\pagestyle{fancy}
\fancyhf{}
\rhead{پروژه پاتوق — سامانه رزرو رستوران}
\lhead{گزارش اسپرینت اول}
\rfoot{\thepage}
\renewcommand{\headrulewidth}{0.4pt}

\begin{document}

% \pagenumbering{persian}  % disabled for hyperref compatibility

%% ─── صفحه عنوان ────────────────────────────────────────────────────────────
\begin{titlepage}
\centering
\vspace*{2cm}
{\LARGE\bfseries پروژه پاتوق (پاتوق)\\[0.5em]}
{\Large\bfseries سامانه مدیریت و رزرو کافه و رستوران\\[2em]}
\rule{\linewidth}{0.5mm}\\[1em]
{\huge\bfseries\color{primaryorange} گزارش اسپرینت اول}\\[0.5em]
{\Large زیرساخت و پایه‌گذاری}\\[1em]
\rule{\linewidth}{0.5mm}\\[2em]
{\large درس: طراحی نرم‌افزار\\[0.5em]}
{\large مدت اسپرینت: ۲ هفته\\[0.5em]}
{\large تاریخ تهیه: تیر ۱۴۰۵}\\[3em]
\vfill
{\normalsize\color{darkgray}
این گزارش پوشش‌دهی اسپرینت اول پروژه پاتوق را شرح می‌دهد که
در آن زیرساخت اصلی بک‌اند و فرانت‌اند، سیستم احراز هویت و
موجودیت‌های اصلی دامنه پیاده‌سازی شدند.}
\end{titlepage}

\tableofcontents
\newpage

%% ─── ۱. هدف اسپرینت ────────────────────────────────────────────────────────
\section{هدف اسپرینت}

هدف اصلی اسپرینت اول، ایجاد پایه‌های محکم معماری سیستم و پیاده‌سازی جریان
احراز هویت کامل است. در پایان این اسپرینت باید کاربر بتواند از طریق \lr{OTP}
یا رمز عبور وارد سیستم شود و توکن‌های دسترسی دریافت کند.

\subsection{تعریف «انجام‌شده» (Definition of Done)}
\begin{itemize}[rightmargin=0.5cm]
  \item کد بدون خطای کامپایل
  \item \lr{API} مستندسازی‌شده در \lr{Swagger}
  \item \lr{Migration}های پایگاه داده اعمال‌شده
  \item صفحات فرانت‌اند با \lr{UI} کارکرد واقعی (نه \lr{mock})
  \item \lr{CORS} صحیح پیکربندی‌شده
\end{itemize}

\newpage

%% ─── ۲. Sprint Backlog ─────────────────────────────────────────────────────
\section{\lr{Sprint Backlog}}

\begin{longtable}{|r|r|r|r|}
\hline
\textbf{آیتم} & \textbf{توضیح} & \textbf{مسئول} & \textbf{وضعیت} \\
\hline
\endhead
راه‌اندازی پروژه بک‌اند & معماری تمیز، لایه‌بندی & بک‌اند & انجام شد \\
\hline
تعریف موجودیت‌های دامنه & کلاس‌های \lr{Domain} & بک‌اند & انجام شد \\
\hline
پیکربندی \lr{EF Core} & \lr{DbContext}، نگاشت‌ها & بک‌اند & انجام شد \\
\hline
پیکربندی \lr{Redis} & \lr{OTP} و کش & بک‌اند & انجام شد \\
\hline
\lr{Docker Compose} & \lr{PostgreSQL} + \lr{Redis} & بک‌اند & انجام شد \\
\hline
سرویس \lr{OTP} & ارسال و تأیید کد & بک‌اند & انجام شد \\
\hline
سرویس \lr{JWT} & صدور توکن دسترسی + تازه‌سازی & بک‌اند & انجام شد \\
\hline
\lr{AuthController} & \lr{Endpoint}های احراز هویت & بک‌اند & انجام شد \\
\hline
\lr{Middleware} خطا & فرمت \lr{RFC 7807} & بک‌اند & انجام شد \\
\hline
راه‌اندازی پروژه فرانت‌اند & \lr{Vite + React + TypeScript} & فرانت‌اند & انجام شد \\
\hline
سیستم مسیریابی & \lr{React Router 7} & فرانت‌اند & انجام شد \\
\hline
صفحه انتخاب نقش & \lr{RoleSelection} & فرانت‌اند & انجام شد \\
\hline
صفحه ورود & \lr{LoginPage} & فرانت‌اند & انجام شد \\
\hline
صفحه تأیید \lr{OTP} & \lr{OTPVerification} & فرانت‌اند & انجام شد \\
\hline
صفحه ثبت‌نام & \lr{RegisterPage} & فرانت‌اند & انجام شد \\
\hline
\lr{AuthContext} & مدیریت وضعیت احراز هویت & فرانت‌اند & انجام شد \\
\hline
\end{longtable}

\newpage

%% ─── ۳. بک‌اند ─────────────────────────────────────────────────────────────
\section{پیاده‌سازی بک‌اند}

\subsection{معماری تمیز (Clean Architecture)}

پروژه بک‌اند بر اساس معماری تمیز با چهار لایه اصلی ساختاردهی شد:

\begin{description}[rightmargin=0.5cm]
  \item[\lr{Patogh.Domain}:] موجودیت‌های دامنه، \lr{Enum}ها و اینترفیس‌های پایه.
    هیچ وابستگی خارجی ندارد.
  \item[\lr{Patogh.Application}:] منطق کسب‌وکار با الگوی \lr{CQRS + MediatR}.
    فقط به \lr{Domain} وابسته است.
  \item[\lr{Patogh.Infrastructure}:] سرویس‌های خارجی: \lr{JWT}، \lr{OTP}، ارسال پیامک،
    آپلود فایل.
  \item[\lr{Patogh.Persistence}:] \lr{DbContext}، \lr{Migration}ها و پیکربندی \lr{EF Core}.
  \item[\lr{Patogh.API}:] کنترلرها، \lr{Middleware}ها و تنظیمات \lr{ASP.NET Core}.
\end{description}

\subsection{موجودیت‌های دامنه}

در این اسپرینت موجودیت‌های اصلی دامنه تعریف شدند:

\begin{longtable}{|r|r|r|}
\hline
\textbf{موجودیت} & \textbf{فیلدهای اصلی} & \textbf{توضیح} \\
\hline
\endhead
\lr{User} & \lr{PhoneNumber, PasswordHash, Role, DisplayName, AvatarUrl} & کاربر سیستم \\
\hline
\lr{Restaurant} & \lr{Name, Description, Location, FoodType, Status, OwnerId} & رستوران \\
\hline
\lr{RestaurantTable} & \lr{TableNumber, Capacity, RestaurantId} & میز رستوران \\
\hline
\lr{MenuItem} & \lr{Name, Description, Price, ImageUrl, RestaurantId} & آیتم منو \\
\hline
\lr{Reservation} & \lr{CustomerId, TableId, Date, StartTime, EndTime, Status} & رزرو \\
\hline
\lr{RefreshToken} & \lr{Token, UserId, ExpiresAt, IsRevoked} & توکن تازه‌سازی \\
\hline
\lr{UserFavorite} & \lr{UserId, RestaurantId} & علاقه‌مندی \\
\hline
\lr{MediaAsset} & \lr{Url, FileName, ContentType} & فایل رسانه‌ای \\
\hline
\end{longtable}

\subsection{پیکربندی پایگاه داده}

\begin{itemize}[rightmargin=0.5cm]
  \item \textbf{\lr{PostgreSQL}:} پایگاه داده اصلی از طریق \lr{Docker Compose} راه‌اندازی شد.
  \item \textbf{\lr{EF Core 8}:} برای \lr{ORM} و مدیریت \lr{Migration}ها.
  \item \textbf{\lr{Redis}:} برای ذخیره‌سازی کدهای \lr{OTP} (با \lr{TTL} 2 دقیقه) و کش.
  \item \textbf{\lr{Docker Compose}:} هر دو سرویس \lr{PostgreSQL} و \lr{Redis} در یک شبکه
    \lr{Docker} اجرا می‌شوند.
\end{itemize}

\subsection{سیستم احراز هویت}

\subsubsection{جریان OTP}
\begin{enumerate}[rightmargin=0.5cm]
  \item کاربر شماره موبایل را ارسال می‌کند (\lr{POST /api/auth/send-otp})
  \item سرور یک کد ۶ رقمی تصادفی تولید و در \lr{Redis} ذخیره می‌کند (کلید: \lr{otp:\{phone\}})
  \item کد از طریق سرویس پیامک ارسال می‌شود
  \item کاربر کد را برای تأیید ارسال می‌کند (\lr{POST /api/auth/verify-otp})
  \item پس از تأیید، کد از \lr{Redis} حذف می‌شود (یکبار مصرف)
  \item توکن دسترسی (\lr{JWT}) و توکن تازه‌سازی صادر می‌شود
\end{enumerate}

\subsubsection{ویژگی‌های امنیتی OTP}
\begin{itemize}[rightmargin=0.5cm]
  \item اعتبار ۲ دقیقه‌ای (\lr{TTL} در \lr{Redis})
  \item یکبار مصرف: پس از تأیید موفق بلافاصله حذف می‌شود
  \item \lr{OTP} در لاگ‌های سیستم ثبت نمی‌شود (امنیت)
  \item محدودیت ۵ درخواست در دقیقه برای جلوگیری از \lr{Brute Force}
\end{itemize}

\subsubsection{ساختار JWT}
\begin{itemize}[rightmargin=0.5cm]
  \item \textbf{توکن دسترسی:} عمر کوتاه (۱۵ دقیقه)، حاوی \lr{UserId، Role، PhoneNumber}
  \item \textbf{توکن تازه‌سازی:} عمر بلند (۷ روز)، ذخیره در پایگاه داده
  \item \lr{Token Rotation}: هر بار استفاده از توکن تازه‌سازی، توکن جدید صادر می‌شود
\end{itemize}

\subsection{Endpoint‌های پیاده‌سازی‌شده}

\begin{longtable}{|r|r|r|r|}
\hline
\textbf{متد} & \textbf{مسیر} & \textbf{توضیح} & \textbf{احراز هویت} \\
\hline
\endhead
\lr{POST} & \lr{/api/auth/send-otp} & ارسال کد \lr{OTP} & بدون نیاز \\
\hline
\lr{POST} & \lr{/api/auth/verify-otp} & تأیید \lr{OTP} و دریافت توکن & بدون نیاز \\
\hline
\lr{POST} & \lr{/api/auth/login} & ورود با رمز عبور & بدون نیاز \\
\hline
\lr{POST} & \lr{/api/auth/register} & ثبت‌نام با رمز عبور & بدون نیاز \\
\hline
\lr{POST} & \lr{/api/auth/refresh} & تجدید توکن دسترسی & بدون نیاز \\
\hline
\lr{POST} & \lr{/api/auth/revoke} & باطل کردن توکن (خروج) & نیاز دارد \\
\hline
\end{longtable}

\newpage

%% ─── ۴. فرانت‌اند ──────────────────────────────────────────────────────────
\section{پیاده‌سازی فرانت‌اند}

\subsection{راه‌اندازی پروژه}

پروژه فرانت‌اند با \lr{Vite 6.3.5 + React 18 + TypeScript} راه‌اندازی شد.

\textbf{پکیج‌های اصلی نصب‌شده:}
\begin{itemize}[rightmargin=0.5cm]
  \item \lr{React Router 7}: مسیریابی سمت کلاینت
  \item \lr{Tailwind CSS v4}: استایل‌دهی با \lr{@tailwindcss/vite}
  \item \lr{shadcn/ui + Radix UI}: کامپوننت‌های \lr{UI}
  \item \lr{React Hook Form}: مدیریت فرم‌ها
  \item \lr{Sonner}: نمایش اعلان‌های \lr{toast}
  \item \lr{Lucide React}: آیکون‌ها
  \item \lr{date-fns}: کار با تاریخ
\end{itemize}

\subsection{سیستم مسیریابی}

مسیریابی پروژه در فایل \texttt{\lr{App.tsx}} تعریف شده است و شامل سه گروه اصلی است:
مسیرهای عمومی، مسیرهای احراز هویت و مسیرهای محافظت‌شده بر اساس نقش.

کامپوننت \texttt{\lr{ProtectedRoute}} از دسترسی کاربران بدون نقش مناسب به صفحات محافظت‌شده
جلوگیری می‌کند.

\subsection{صفحات پیاده‌سازی‌شده}

\subsubsection{صفحه انتخاب نقش (RoleSelection)}
اولین صفحه در جریان ورود، دو دکمه برای انتخاب نقش مشتری یا مدیر رستوران نمایش می‌دهد
و کاربر را به صفحه ورود مناسب هدایت می‌کند.

\subsubsection{صفحه ورود (LoginPage)}
\begin{itemize}[rightmargin=0.5cm]
  \item دو حالت: ورود با \lr{OTP} (پیش‌فرض) و ورود با رمز عبور
  \item اعتبارسنجی فرمت شماره موبایل ایرانی
  \item ارسال درخواست \lr{OTP} و هدایت به صفحه تأیید
  \item نمایش پیام‌های خطای فارسی
\end{itemize}

\subsubsection{صفحه تأیید OTP (OTPVerification)}
\begin{itemize}[rightmargin=0.5cm]
  \item اینپوت ۶ خانه‌ای برای کد \lr{OTP}
  \item تایمر شمارش معکوس ۲ دقیقه‌ای
  \item امکان درخواست ارسال مجدد کد
  \item هدایت خودکار پس از تأیید موفق بر اساس نقش
\end{itemize}

\subsubsection{صفحه ثبت‌نام (RegisterPage)}
فرم ثبت‌نام با شماره موبایل، رمز عبور و تأیید رمز عبور. اعتبارسنجی تمام فیلدها
با \lr{React Hook Form}.

\subsection{مدیریت وضعیت احراز هویت}

\lr{AuthContext} برای مدیریت وضعیت ورود کاربر در تمام صفحات ایجاد شد:

\begin{itemize}[rightmargin=0.5cm]
  \item ذخیره توکن دسترسی در \lr{localStorage}
  \item ذخیره اطلاعات کاربر (نقش، شماره، نام)
  \item تابع خروج از سیستم
  \item بارگذاری خودکار وضعیت از \lr{localStorage} پس از رفرش صفحه
\end{itemize}

\newpage

%% ─── ۵. چالش‌ها ─────────────────────────────────────────────────────────────
\section{چالش‌ها و راه‌حل‌ها}

\subsection{چالش ۱: خطای BCrypt با هش خالی}
\begin{description}[rightmargin=0.5cm]
  \item[مشکل:] هنگام ثبت‌نام کاربر بدون رمز عبور (ورود با \lr{OTP})، فراخوانی
    \lr{BCrypt.Verify} با رشته خالی خطا می‌داد.
  \item[راه‌حل:] افزودن بررسی \texttt{\lr{string.IsNullOrEmpty(PasswordHash)}} قبل از
    فراخوانی \lr{BCrypt}. کاربران \lr{OTP}-only مقدار \lr{PasswordHash} خالی دارند و
    ورود با رمز عبور برای آن‌ها امکان‌پذیر نیست.
\end{description}

\subsection{چالش ۲: پیکربندی CORS}
\begin{description}[rightmargin=0.5cm]
  \item[مشکل:] فرانت‌اند روی پورت \lr{5173} اجرا می‌شد اما \lr{CORS} سرور این پورت را
    مجاز نمی‌دانست.
  \item[راه‌حل:] پیکربندی \lr{CORS} با افزودن \lr{http://localhost:5173} به لیست
    \lr{AllowedOrigins} در \lr{appsettings.json}. همچنین \lr{AllowCredentials()} برای
    ارسال کوکی‌های توکن تازه‌سازی فعال شد.
\end{description}

\subsection{چالش ۳: مدیریت وضعیت IsNewUser}
\begin{description}[rightmargin=0.5cm]
  \item[مشکل:] تشخیص اینکه کاربر پس از تأیید \lr{OTP} باید به صفحه تکمیل پروفایل هدایت
    شود یا به داشبورد.
  \item[راه‌حل:] فیلد \lr{HasIncompleteProfile} در پاسخ \lr{verify-otp} اضافه شد.
    فرانت‌اند بر اساس این فیلد تصمیم‌گیری می‌کند.
\end{description}

\newpage

%% ─── ۶. بازبینی و گذر ───────────────────────────────────────────────────────
\section{بازبینی و گذر اسپرینت}

\subsection{Sprint Review}
\begin{itemize}[rightmargin=0.5cm]
  \item تمامی \lr{Endpoint}های احراز هویت پیاده‌سازی و تست شدند
  \item جریان کامل \lr{OTP} از ارسال تا دریافت توکن کار می‌کند
  \item صفحات فرانت‌اند به \lr{API} واقعی وصل شدند
  \item \lr{Docker Compose} برای محیط توسعه پیکربندی شد
\end{itemize}

\subsection{Sprint Retrospective}

\textbf{چه چیزی خوب بود:}
\begin{itemize}[rightmargin=0.5cm]
  \item معماری تمیز از ابتدا پیاده‌سازی شد که نگهداری کد را ساده کرد
  \item استفاده از \lr{MediatR} کدهای کنترلر را کوچک و تمیز نگه داشت
  \item \lr{Docker Compose} راه‌اندازی محیط توسعه را سریع کرد
\end{itemize}

\textbf{چه چیزی می‌توانست بهتر باشد:}
\begin{itemize}[rightmargin=0.5cm]
  \item پیکربندی \lr{CORS} باید زودتر انجام می‌شد تا تأخیر در اتصال فرانت به بک‌اند کاهش می‌یافت
  \item باید از ابتدا برای \lr{HasIncompleteProfile} برنامه‌ریزی می‌شد
\end{itemize}

\textbf{اقدامات بهبود برای اسپرینت بعدی:}
\begin{itemize}[rightmargin=0.5cm]
  \item تعریف مشخص‌تر \lr{API contract} قبل از شروع پیاده‌سازی
  \item نوشتن \lr{Integration Test} برای \lr{Endpoint}های اصلی
\end{itemize}

\end{document}
